\section{Conclusions and outlook}
\label{sec:conclusions}

We have studied how the equilibrium structure, linear stability, and local
periodic dynamics of an infinitesimal fifth body change as a general
non-symmetric four-primary central configuration is continuously deformed.
Rather than comparing isolated primary configurations, the analysis follows
the admissible central-configuration set itself and resolves the restricted
equilibria, their bifurcations, and selected periodic-orbit families along
that set.

Three conclusions emerge.

First, on the investigated slice
\[
(a,b)=(0.9,-0.1),
\]
the admissible central-configuration set is numerically separated into two
distinct components containing the Case-1 and Case-2 reference
configurations. The five- and nine-equilibrium cases therefore do not lie
on a single fixed-$(a,b)$ continuation family. This establishes that
equilibrium multiplicity must be interpreted together with the component
structure of the underlying primary configuration: different sampled
counts need not be connected by a bifurcation sequence on the same
admissible component. The result is specific to the investigated slice and
does not establish disconnection in the full four-parameter admissible set.

Second, the restricted equilibrium structure along these components is
organized by seven generic saddle-node folds. The corresponding interior
multiplicity sequences are
\[
5\rightarrow7\rightarrow5
\]
on the Case-1 component and
\[
7\rightarrow9\rightarrow11\rightarrow9\rightarrow7\rightarrow5
\]
on the Case-2 component. The Case-1 arc additionally exhibits a
mass-vanishing boundary degeneration as $\rho_2\rightarrow0$, which is
distinct from the interior folds. The eleven-equilibrium interval is
structurally significant because it contains a branch exchange: the pair
$B_a$ created at the $9\rightarrow11$ fold is not the pair $B_b$
destroyed at the following $11\rightarrow9$ fold. Thus two configurations
can have the same equilibrium count while containing different sets of
equilibrium branches. The scalar multiplicity is therefore only a coarse
description of the underlying equilibrium structure.

Third, this equilibrium organization is inherited by the local dynamics.
Both admissible components contain reliably resolved intervals of
spectrally stable equilibria, and the Case-2 component supports two
distinct spectrally stable equilibrium branches over overlapping parameter
ranges. From selected simple centre modes we computed ten Lyapunov
periodic-orbit families comprising 198 corrected solutions. The families
continued from the selected doubly-elliptic parents remain Floquet stable
over the computed intervals, whereas the families continued from selected
saddle$\times$centre parents possess a real reciprocal Floquet pair and are
real unstable.

Low-order centre-frequency ratios encountered along the stable branches are
near-commensurabilities rather than exact resonances at the parent
equilibria. Families~7--8 sample the Case-1 region near $3{:}2$, while
families~9--10 sample the distinct Case-2 region near $5{:}2$. Family~10
approaches the period-doubling boundary extremely closely,
\[
\min |\nu+2|\simeq10^{-7},
\]
but an ultra-fine amplitude sweep resolves a turn-back rather than a
crossing. No period-doubling bifurcation is therefore detected on the
computed portion of that family. More generally, no crossing of the
selected low-order multiplier conditions is found along the ten continued
families; this does not exclude secondary periodic solutions elsewhere in
parameter or phase space.

The equilibrium structure also determines a hierarchy of critical Jacobi
levels. A representative calculation in the eleven-equilibrium regime
shows a resolution-robust net change in Hill-region connectivity from one
component to two across a small cluster of three nearby equilibrium
critical levels. Representative Floquet-stable periodic orbits remain
inside the corresponding allowed regions at their own Jacobi constants.
This provides a concrete geometric link between the equilibrium
bifurcations, the critical energy levels, the accessible Hill geometry, and
the periodic dynamics without implying that every critical level produces
a change in component count.

Taken together, the results support a broader interpretation: in a
non-symmetric restricted problem, the topology of the admissible
primary-configuration set can itself act as an organizing structure for the
restricted dynamics. On the present slice it determines which equilibrium
branches can coexist and which bifurcation sequences are accessible; those
branches in turn determine the available centre modes, critical Jacobi
levels, and local periodic families. Whether this organizing role persists
in the full parameter space is a natural question for further study.

Several extensions follow directly from the present analysis. The most
important is continuation in the complete four-parameter admissible
central-configuration set, which would determine whether the two components
identified here remain disconnected when $a$ and $b$ are also allowed to
vary. A second direction is pseudo-arclength continuation of the
periodic-orbit families themselves, together with targeted continuation of
secondary branches near the observed low-order commensurabilities and the
near-period-doubling point. A third is a systematic Hill-region
classification across individual critical Jacobi levels rather than the
representative connectivity calculation performed here. Finally, comparison
with symmetric restricted five-body models at matched or otherwise
controlled mass configurations would help separate effects caused by mass
distribution from those caused specifically by the geometry and topology
of the admissible central-configuration set.
